\subsection{Exercise 5 - Lab 3}

Yêu cầu: Trong thực tế, các nút nhấn cơ học khi thao tác luôn xảy ra hiện tượng dội tín hiệu trong khoảng $10\text{ms} \rightarrow 20\text{ms}$ đầu tiên, làm vi điều khiển nhận diện nhầm một lần nhấn thành nhiều lần nhấn liên tục.

Để giải quyết triệt để hiện tượng này, bài tập 5 triển khai thuật toán lấy mẫu 3 lần liên tiếp bên trong ngắt Timer 2 với chu kỳ $10\text{ms}$:
\begin{enumerate}
    \item Cập nhật 3 thanh ghi lấy mẫu: \texttt{KeyReg0}, \texttt{KeyReg1}, \texttt{KeyReg2}.
    \item Chỉ xác nhận trạng thái nút nhấn khi cả 3 thanh ghi thu được cùng một giá trị ổn định trong ít nhất $30\text{ms}$ ($3 \times 10\text{ms}$).
    \item Sử dụng thanh ghi \texttt{KeyReg3} để lưu trạng thái ổn định và xác định thời điểm bắt đầu nhấn (\textit{Falling edge detection}) để kích hoạt cờ \texttt{button\_flag}.
\end{enumerate}

\textbf{Mã nguồn Cài đặt}
\begin{lstlisting}[language=C, caption={File input reading.h}]
#ifndef INC_INPUT_READING_H_
#define INC_INPUT_READING_H_

#include"main.h"
#include "timer.h"

void init_buffer();
void button_reading(void);
unsigned char is_button_pressed(unsigned char index);
unsigned char is_button_pressed_1s(unsigned char index);
unsigned char is_button_held(unsigned char index);
void reset_flagForButtonHold(unsigned char index);

#endif /* INC_INPUT_READING_H_ */
\end{lstlisting}

\begin{lstlisting}[language=C, caption={Cài đặt lọc nhiễu dội phím trong file input reading.c}]
static GPIO_PinState buttonBuffer[NO_OF_BUTTONS] = {BUTTON_IS_RELEASED, BUTTON_IS_RELEASED, BUTTON_IS_RELEASED};
static GPIO_PinState debounceButtonBuffer1[NO_OF_BUTTONS] = {BUTTON_IS_RELEASED, BUTTON_IS_RELEASED, BUTTON_IS_RELEASED};
static GPIO_PinState debounceButtonBuffer2[NO_OF_BUTTONS] = {BUTTON_IS_RELEASED, BUTTON_IS_RELEASED, BUTTON_IS_RELEASED};
static GPIO_PinState debounceButtonBuffer3[NO_OF_BUTTONS] = {BUTTON_IS_RELEASED, BUTTON_IS_RELEASED, BUTTON_IS_RELEASED};

static uint8_t flagForButtonPress1s[NO_OF_BUTTONS] = {0, 0, 0};
static uint16_t counterForButtonPress1s[NO_OF_BUTTONS] = {0, 0, 0};
static uint8_t flagForButtonHold[NO_OF_BUTTONS] = {0, 0, 0};
static uint16_t counterForButtonHold[NO_OF_BUTTONS] = {0, 0, 0};
\end{lstlisting}

\begin{lstlisting}[language=C, caption={Hàm xử lý phím trong file input reading.c}]
void button_reading(void) {
    for (int i = 0; i < NO_OF_BUTTONS; i++) {
        debounceButtonBuffer3[i] = debounceButtonBuffer2[i];
        debounceButtonBuffer2[i] = debounceButtonBuffer1[i];


        debounceButtonBuffer1[i] = HAL_GPIO_ReadPin(button_ports[i], button_pins[i]);

        if (debounceButtonBuffer1[i] == debounceButtonBuffer2[i] && debounceButtonBuffer3[i] == debounceButtonBuffer2[i])
            buttonBuffer[i] = debounceButtonBuffer1[i];

        if (buttonBuffer[i] == BUTTON_IS_PRESSED) {
            if (counterForButtonPress1s[i] < DURATION_FOR_AUTO_INCREASING) {
                counterForButtonPress1s[i]++;
            } else {
                flagForButtonPress1s[i] = 1;
                if (counterForButtonHold[i] < DURATION_FOR_HOLD) {
                    counterForButtonHold[i]++;
                    if (counterForButtonHold[i] >= DURATION_FOR_HOLD) {
                        counterForButtonHold[i] = 0;
                        flagForButtonHold[i] = 1;
                    }
                }
            }
        } else {
            counterForButtonPress1s[i] = 0;
            flagForButtonPress1s[i] = 0;
            counterForButtonHold[i] = 0;
            flagForButtonHold[i] = 0;
        }
    }
}
unsigned char is_button_pressed(uint8_t index){
    if(index >= NO_OF_BUTTONS) return 0;
    return (buttonBuffer[index] == BUTTON_IS_PRESSED);
}

unsigned char is_button_pressed_1s(unsigned char index){
    if(index >= NO_OF_BUTTONS) return 0xff;
    return (flagForButtonPress1s[index] == 1);
}

unsigned char is_button_held(unsigned char index) {
    if (index >= NO_OF_BUTTONS) return 0;
    return (flagForButtonHold[index] == 1);
}
\end{lstlisting}

